% ============================================================
%  Publication 3 — Bahá'í Studies
%  "Mahabbat and the Two Wings: Quantum Physics as
%   Structural Convergence with Bahá'í Relational Ontology"
%
%  Target: Journal of Bahá'í Studies · Lights of Irfan ·
%          Association for Bahá'í Studies Conference
%  Compiler: pdflatex (Overleaf)
%  No mathematics required for validation of this paper.
%  Dr. Joshua Adams | 2026
% ============================================================

\documentclass[12pt, a4paper]{article}

% ── Encoding & fonts ──────────────────────────────────────
\usepackage[T1]{fontenc}
\usepackage[utf8]{inputenc}

% ── Page geometry ─────────────────────────────────────────
\usepackage[margin=1in]{geometry}
\usepackage{setspace}
\doublespacing

% ── Tables & figures ──────────────────────────────────────
\usepackage{graphicx}
\usepackage{booktabs}
\usepackage{tabularx}
\graphicspath{{figures/}}

% ── Arabic / Persian text support ─────────────────────────
% (mahabbat rendered in Latin transliteration; no arabtex needed)

% ── Bibliography (author-year for humanities journals) ────
\usepackage[authoryear, round, sort]{natbib}
\bibliographystyle{apalike}

% ── Hyperlinks ────────────────────────────────────────────
\usepackage{hyperref}
\hypersetup{colorlinks=true, linkcolor=blue, citecolor=blue, urlcolor=blue}

% ── ORCID ─────────────────────────────────────────────────
\usepackage{orcidlink}

% ── Bahá'í typesetting helpers ────────────────────────────
\usepackage{xspace}
\newcommand{\bahai}{Bah\'a\textquoteright\'{\i}\xspace}
\newcommand{\Bahai}{Bah\'a\textquoteright\'{\i}\xspace}
\newcommand{\Bahaullah}{Bah\'a\textquoteright u\textquoteright ll\'ah\xspace}
\newcommand{\AbdulBaha}{\textquoteleft Abdu\textquoteright l-Bah\'a\xspace}
\newcommand{\mahabbat}{\textit{mahabbat}\xspace}
\newcommand{\mawaddat}{\textit{mawaddat}\xspace}

% =============================================================
\begin{document}
% =============================================================

\title{Mahabbat and the Two Wings: Quantum Physics as Structural\\
       Convergence with Bahá'í Relational Ontology}

\author{Dr. Joshua Adams \orcidlink{0000-0002-7185-9125}\\
        \small Independent Researcher\\
        \small \texttt{contact@joshuaadams.dev}}

\date{\today}

\maketitle

\begin{abstract}
\AbdulBaha taught that religion and science are the ``two wings'' of one truth. This paper tests that principle at the level of structure. It argues that the \Bahai metaphysics of \mahabbat---relational love as the foundation of existence---shares the same logical architecture as relational quantum mechanics and Whitehead's process philosophy: a world constituted by relationship rather than substance, sustained by love. The Writings articulate this structure decades before Whitehead and a century before physics confirmed it experimentally. This is offered not as proof of supernatural authority but as documented independent convergence---the two-wings principle confirmed in structure, not metaphor.

\end{abstract}

\newpage
\tableofcontents
\newpage

% ── SECTION 1: INTRODUCTION — THE TWO-WINGS TEST (~700 words) ───────────────

\section{Introduction: The Two-Wings Test}
\label{sec:intro}

% Key question: what would genuine confirmation of the two-wings principle
% look like? Not shared conclusions — but shared structure.
% Introduce structural isomorphism vs. analogy.
% Preview six convergences and chronological sequence.

\AbdulBaha taught that science and religion are ``like the two wings upon which
the bird of human intelligence can soar into the heaven of truth.''
The principle is foundational to \bahai teaching and frequently invoked.
But what would it mean to confirm it?

Previous \bahai-science dialogue has often worked by identifying resonances:
quantum uncertainty sounds like the unknowability of God; entanglement sounds
like the interconnection of all things; the Big Bang sounds like creation.
These resonances are suggestive.
But they are analogical, not structural.
Analogy shows that two things look similar; structural isomorphism shows that
they have the same logical form.

This paper asks whether the two-wings principle can be confirmed at the level
of structure --- whether the \bahai metaphysical framework and the framework
that modern physics has empirically arrived at share the same logical
architecture, arrived at by independent means.

The answer I document is yes.
On my reading, the \bahai metaphysics of \mahabbat --- relational love as the creative
foundation of existence at every degree --- shares the same logical structure
as relational quantum mechanics~\citep{Rovelli1996} and process
ontology~\citep{Whitehead1929}: existence constituted by relationship,
properties emerging through relational interaction, the force sustaining
relationship identified as love.

The paper proceeds in five movements.
\S\ref{sec:mahabbat} establishes, on my reading, the \bahai metaphysics of \mahabbat through
close reading of three core passages and their secondary corroborations.
\S\ref{sec:convergences} documents six structural correspondences between that
metaphysics and the framework of contemporary physics: relational ontology,
nonlocal interconnection, observer-dependence and the hierarchy of
comprehension, decoherence and the dissolution of forms, the holographic and
implicate orders, and motion as the character of existence.
\S\ref{sec:chronology} places the convergences on a timeline that begins in
the \bahai revelation of the late nineteenth century, passes through
\citeauthor{Whitehead1929}'s philosophical articulation in 1929, and arrives
at the experimental confirmation that earned the 2022 Nobel Prize in Physics.
\S\ref{sec:confirmation} addresses the methodological question raised by the
convergence and distinguishes it from any claim of borrowing or appropriation.
\S\ref{sec:soul} considers one implication for \bahai scholarship: the
structural status of the non-composite soul within a relational framework.
\S\ref{sec:conclusion} returns to the two-wings principle and what its
confirmation at the level of structure --- as distinct from the level of
metaphor --- contributes to academic \bahai-science dialogue.

% ── SECTION 2: THE BAHÁ'Í METAPHYSICS OF MAHABBAT — TEXTUAL ANALYSIS (~1,800 words)

\section{The \bahai Metaphysics of \textit{Mahabbat}: Textual Analysis}
\label{sec:mahabbat}

% Close reading of three core passages + secondary passages.
% Distinguish mahabbat (metaphysical) from emotional love.
% Establish the ontological claim: existence IS relational affinity.

My argument rests on a specific reading of three passages from
the \bahai writings.
The reading is structural: I read each passage as a precise ontological
claim about the constitution of reality, not as a poetic or devotional figure.
The passages were not delivered as physics, and they do not pretend to be physics.
But they are propositions, and they have a logical form, and that form is
what concerns me here.

\subsection{The Three Core Passages}

\paragraph{Existence as relational affinity.}
\AbdulBaha writes:
\begin{quote}
``All created things are expressions of the affinity and cohesion of
elementary substances, and nonexistence is the absence of their
attraction and agreement. Everything partakes of this nature and is
subject to this principle, for the creative foundation in all its
degrees and kingdoms is an expression or outcome of love.''
\citep[Talk 48]{PUP}
\end{quote}

Three claims are stated here in succession, and each is structural.
First, existence is defined positively as the ongoing affinity and cohesion of
elementary substances --- not as a property those substances possess, but as the
relational event in which they participate.
Second, nonexistence is defined negatively as the absence of that affinity:
not as a different state of being, but as the cessation of the relational event
itself.
Third, the principle is declared universal: ``everything partakes of this
nature \ldots in all its degrees and kingdoms.''
The scope is not metaphorical extension from one level to another.
It is the assertion of a single relational architecture operating from the
mineral kingdom to the human soul to the divine.
And it is named: the principle is ``an expression or outcome of love.''
What constitutes existence at every degree is identified, in plain language,
as \mahabbat.

This is not a poetic gloss on an otherwise substantialist metaphysics.
The passage refuses, in three successive moves, to allow the reader to locate
existence in any substrate beneath the relation.
There are no elementary substances first, and affinity second.
Affinity \emph{is} how the elementary substances are; absent the affinity, there
are no substances to speak of, only nonexistence.
The Bell-theorem reader will recognise the structure: properties do not inhere
in isolated systems and then enter into relations; properties are constituted
through the relations themselves~\citep{Bell1964,Rovelli1996}.
I read \citeauthor{Rovelli1996} as articulating the same structural principle in 1996.
On my reading, \AbdulBaha had articulated it in 1912.

\paragraph{Motion and process as the nature of reality.}
\begin{quote}
``Creation is the expression of motion. Motion is life. A moving object
is a living object, whereas that which is motionless and inert is as dead.
All created forms are progressive in their planes, or kingdoms of
existence, under the stimulus of the power or spirit of life. The
universal energy is dynamic. Nothing is stationary in the material
world of outer phenomena or in the inner world of intellect and
consciousness.''
\citep[Talk 52]{PUP}
\end{quote}

The passage moves from creation to motion to life in three sentences and does
not return.
Reality is not described as a set of things that happen to undergo motion.
It is described as motion itself --- as a process whose essential character is
the continuous becoming that any moment of inertness would interrupt.
The grammar is direct: creation \emph{is} motion; motion \emph{is} life.
The identifications are not similes.
They are equations of essence.

Physics would establish, over the half-century that followed, that the
prohibition of perfect rest is not a metaphysical preference but a theorem.
\citeauthor{Heisenberg1927}'s uncertainty principle forbids any system from
occupying a state of zero motion: the vacuum itself fluctuates, and the lowest
energy state of any physical system retains an irreducible residual
agitation~\citep{Heisenberg1927}.
The passage and the theorem agree on the conclusion that reality is, at the
deepest level of physical description, dynamic through and through.
The passage states it directly; the theorem derives it from the commutation
relations of quantum observables.
That two such different methods agree on the same conclusion is the kind of
convergence that I document in this paper.

\paragraph{Creation as relational event.}
\begin{quote}
``The world of existence came into being through the heat generated
from the interaction between the active force and that which is its recipient.
These two are the same, yet they are different.''
\citep[p.~140]{TabletsBahaullah}
\end{quote}

Two structural claims are made here.
The first is that creation is an event of relational interaction --- the
``heat generated'' is the product of the relation, not of either relatum in
isolation.
The second is the paradox: ``these two are the same, yet they are different.''
Active and recipient are distinguished sufficiently for the relation to obtain,
and yet not so distinguished that they fall apart into independent entities.
On my reading, the same-yet-different structure is analogous to the structure of a quantum
entangled pair: two correlated systems whose joint state is a single
indivisible reality, and yet which can be addressed under partition as two.
There is no third option in either case.
There are not two prior things that subsequently relate, nor one prior thing
that subsequently divides.
There is the relational unity, and within it, the differentiation.

The Tablet of Wisdom from which the passage derives is among
\Bahaullah's most cosmologically explicit writings~\citep{TabletsBahaullah}.
The reading I offer here treats it not as figurative theology but as a direct
description of the relational ontology that physics would later confirm by
experiment.
The reader is free, on philosophical or theological grounds, to prefer a
figurative reading.
What is not available is the claim that no structural reading is possible.
The structure is in the text.

\subsection{Secondary Passages}

Three further passages corroborate the reading.

\Bahaullah writes in the \textit{Hidden Words}:
``O Son of Man! I loved thy creation, hence I created thee.
Wherefore, do thou love Me, that I may name thy name and fill thy soul
with the spirit of life''~\citep[Arabic 4]{HiddenWords}.
The passage locates the originating event of creation in love and identifies
the answering motion of the creature as love returned.
On my reading, it states that the originating affinity is reciprocal and
that the reciprocity is constitutive:
the creature is not first made and then loved; the love is the making.

\AbdulBaha returns, in a complementary passage, to the consequences of
the absence of affinity:
``As long as there is affinity and cohesion among these constituent
elements, strength and life are manifest; but when dissension and
repulsion arise among them, disintegration follows.''~\citep[Talk 41]{PUP}
The passage is the converse of the Talk 48 passage analysed above.
If existence is the ongoing affinity of constituents, then the cessation of
that affinity is the dissolution of the form.
The passage does not invoke an external cause of destruction.
Disintegration is not done to a system; it is what a system becomes when the
affinity sustaining it fails.
This, I argue, has the structural shape of decoherence, to which I return in
\S\ref{sec:convergences}.

Finally, the \textit{Tablet of Wisdom} extends the principle from individual
creation events to the cosmos itself~\citep{TabletsBahaullah}.
\Bahaullah declares the unity of the world and the relational character of
its constituents in terms that, on my reading, can be read alongside the holographic and implicate
orders that physics would propose a century later~\citep{Bohm1980,Susskind1995}.
I do not press the comparison here.
What the passage establishes for present purposes is that the relational
architecture of \mahabbat is offered not as a feature of human experience but
as the structure of created reality as such.

\subsection{\textit{Mahabbat} vs.\ Emotional Love}

The argument depends on a distinction that the Arabic vocabulary of the
\bahai writings makes available but English flattens.
\mahabbat is one of several Arabic terms for love, and it is the term most
consistently deployed in the writings when the relational foundation of
existence is at issue.
It is not interchangeable with \textit{hubb} (general affection),
\mawaddat (affectionate kindness, often interpersonal), or
\textit{\textquoteleft ishq} (the ardent, consuming love of the Sufi mystical
tradition, which \Bahaullah deploys deliberately in \textit{The Seven Valleys}).
The distinctions are not pedantic.
\textit{Mahabbat} carries the semantic weight of cordial inclination as a
sustained relational disposition, and it is in this register that the
writings invoke it when speaking of the creative foundation.

The passage from Talk 48, on a careful reading, does not say that the
creative foundation feels affectionately toward what it creates.
It says that the creative foundation \emph{is} the relational affinity by
which constituents cohere.
What I call emotional love here --- the felt experience of attraction,
care, and inclination toward another --- is on my reading the most
conscious and most intensified instance of the same relational force that
constitutes the existence of an atom.
The relation between the metaphysical and the experiential senses is
therefore not analogy.
It is identity of structure at different degrees of expression.
The atom does not feel love; but what holds the atom together is the same
relational affinity that, at the level of human consciousness, registers
as feeling.

This collapses, or at least relocates, a familiar problem.
On a substantialist metaphysics, the felt experience of love is a subjective
overlay on an indifferent physical substrate, and its apparent depth is
either a useful evolutionary fiction or an explanatory remainder.
On the relational metaphysics articulated by the \bahai writings, the felt
experience of love is the most direct phenomenological access we have to the
relational structure that physics has independently confirmed at the level of
elementary systems.
The experience is not a projection onto an indifferent world.
It is the world becoming explicitly conscious of what it has always been doing.

% ── SECTION 3: THE SIX CONVERGENCES FROM BAHÁ'Í TEXTS (~2,000 words) ────────

\section{The Six Convergences}
\label{sec:convergences}

% For each convergence: \bahai textual evidence + accessible physics summary.
% Audience: \bahai scholars. Physics stated accessibly without full formalism.

I turn now to the convergences themselves.
I argue that six structural correspondences obtain between the \bahai metaphysics of
\mahabbat and the framework that physics has empirically established.
Each is presented in three parts: the \bahai claim with its textual source,
the physics with its primary citations, and a brief statement of the
structural identity between them.
The physics is stated accessibly; no formalism is required for the argument
to be assessed.

\subsection{Convergence 1: Relational Ontology}

\paragraph{The \bahai claim.}
``Spirit, mind, soul, and the powers of sight and hearing are but one single reality which hath manifold expressions owing to the diversity of its instruments''~\citep[para.~35]{SummonsLordHosts}.

\paragraph{The physics.}
Bell's theorem proves that quantum properties are not intrinsic to isolated
systems --- they emerge through relational interaction~\citep{Bell1964,Rovelli1996}.
On my reading, reality, at its most fundamental level, is irreducibly relational.

\subsection{Convergence 2: Nonlocal Interconnection}

\paragraph{The \bahai claim.}
``All created things are connected one to another by a linkage complete and
perfect''~\citep[\S 21]{SelectionsAbdul}.

\paragraph{The physics.}
Entangled quantum particles maintain nonlocal correlations across any distance~\citep{Aspect1982,Zeilinger2023Nobel}.
This is not action at a distance; it is the signature of a relational bond
that persists regardless of spatial separation.

\subsection{Convergence 3: Hierarchy of Comprehension and Observer-Dependence}

\paragraph{The \bahai claim.}
% VERIFIED (this session) against SAQ Part 4 "On the Origin, Powers, and
% Conditions of Man" at bahai.org: the graded comprehension teaching is
% explicit there (e.g. the animal "cannot comprehend or conceive that which
% lies beyond it"). Part-level citation confirmed correct.
\AbdulBaha articulates a graded hierarchy of comprehension extending from
mineral to vegetable to animal to human, with each kingdom able to
apprehend the kingdoms beneath it but not the kingdom above it~\citep[Part 4]{SAQ}.
What is fully real for one degree of comprehension is partially or wholly
inaccessible to another.

\paragraph{The physics.}
Relational quantum mechanics establishes that the description of a system is
always relative to some other system that interacts with it~\citep{Rovelli1996}.
What becomes definite, and to whom, depends on the relational context.
Two observers may consistently assign different quantum states to the same
system, not because one is mistaken but because each description is correct
relative to its own reference frame.

\paragraph{The convergence.}
I read both frameworks as abandoning the view that reality presents one and the same face
to every observer.
What an observer can register, and at what degree of articulation, depends on
the observer's own relational and structural capacities.
On my reading, the \bahai hierarchy and the relational interpretation agree on the structural
point: there is no view from nowhere.

\subsection{Convergence 4: Decoherence and Decomposition}

\paragraph{The \bahai claim.}
``As long as there is affinity and cohesion among these constituent
elements, strength and life are manifest; but when dissension and
repulsion arise among them, disintegration follows.''~\citep[Talk 41]{PUP}
\AbdulBaha adds, however, that ``total annihilation is an impossibility'';
what disintegrates is the form, not the constituents in their underlying
relational being~\citep[Talk 38]{PUP}.

\paragraph{The physics.}
Quantum decoherence is the process by which a system loses its capacity to
exhibit interference, as its internal degrees of freedom become entangled with
the degrees of freedom of its environment~\citep{Zurek2003}.
The coherent superposition is not destroyed; it spreads into correlations the
local observer can no longer recover.
The information is conserved; what disperses is not the system's identity but its physical, localized form---the coherent structure that had made it detectable as a bounded object.

\paragraph{The convergence.}
I read the structural shape as identical.
A coherent form is sustained by an ongoing affinity among its constituents;
when that affinity dissipates into the environment, the form perishes; what
underlies the form is not annihilated but rendered inaccessible to local
description.
On my reading, the \bahai principle that nothing is annihilated even when forms dissolve
agrees, term for term, with the conservation of quantum information under
decoherence.

\subsection{Convergence 5: Holographic and Implicate Order}

\paragraph{The \bahai claim.}
``Within thee have I placed the essence of My light''~\citep[Arabic 12]{HiddenWords}.
The \bahai writings repeatedly affirm that the whole of reality is in some
manner present in each of its parts, and that the human soul is the locus in
which the universe becomes interior to itself.

\paragraph{The physics.}
Bohm's implicate order proposes that the explicate world we perceive is
the unfolding of a deeper, enfolded order in which every region carries
information about the whole~\citep{Bohm1980}.
The holographic principle, developed independently in high-energy physics,
establishes that the information content of any spatial volume can be encoded
on its boundary surface~\citep{tHooft1993,Susskind1995}.
In both proposals, the whole is implicit in the part.

\paragraph{The convergence.}
This convergence is the most analogical of the six and I present it with
appropriate caution.
What the \bahai passages assert is the structural availability of the whole
within the human interior, and what the physics establishes is a corresponding
structural availability of global information at every local scale.
Whether the two are describing the same reality from different vantages, or
two related but distinct features of a relational cosmos, is a question I
leave open.
The structural correspondence, on my reading, is real.

\subsection{Convergence 6: Motion as Life}

\paragraph{The \bahai claim.}
``Creation is the expression of motion. Motion is life. A moving object is a
living object, whereas that which is motionless and inert is as dead.''
\citep[Talk 52]{PUP}
The passage was analysed in \S\ref{sec:mahabbat}; I restate it here to mark
its place in the convergence structure.

\paragraph{The physics.}
\citeauthor{Heisenberg1927}'s uncertainty principle forbids any quantum system
from occupying a state of zero motion: a system localised perfectly in position
would have undefined momentum and infinite energy~\citep{Heisenberg1927}.
The quantum vacuum --- the lowest energy state available to a field --- is not
empty but seethes with persistent fluctuations.
There is no rest at the bottom of physical description.

\paragraph{The convergence.}
The most direct of the six.
I read the \bahai claim that creation \emph{is} motion and that inertness is the
character of nonexistence as making the same structural claim that physics derives
from the commutation relations of quantum observables: reality is dynamic
through and through, at every scale.
The two statements arrive by entirely different routes at the same conclusion.

% ── SECTION 4: THE CHRONOLOGICAL SEQUENCE (~700 words) ──────────────────────

\section{The Chronological Sequence: What It Shows}
\label{sec:chronology}

% Present convergence table in narrative form.
% What the sequence shows and what it does not prove.

The convergence I have documented has a chronological dimension
that is part of my argument.
A finding of structural identity between two traditions might, in principle,
be explained by influence: one tradition learns the structure from the other
and reformulates it in its own vocabulary.
The chronology shown in Table~\ref{tab:chronology} forecloses that explanation.
The \bahai articulations precede the philosophical and physical articulations
by margins that make influence impossible in either direction.
(The characterizations of the \bahai writings in the table reflect my analytical reading;
the texts themselves did not employ the vocabulary of process ontology or relational QM.)

\begin{table}[h]
\centering
\caption{Chronological sequence of convergent findings.}
\label{tab:chronology}
\begin{tabularx}{\textwidth}{llX}
\toprule
\textbf{When} & \textbf{By Whom / How} & \textbf{What Was Established} \\
\midrule
1850s--1892  & \Bahaullah (revelation)
             & Process ontological claims: reality as motion, creation as
               relational event, existence as relational affinity \\
1892--1921   & \AbdulBaha (teaching)
             & Love as creative foundation; hierarchy of comprehension;
               soul as non-composite \\
1929         & Whitehead (philosophical analysis)
             & \textit{Process and Reality}: actual occasions, prehension,
               eros --- process ontology named philosophically \\
1964         & Bell (mathematics)
             & Bell's theorem: no local hidden-variable theory reproduces quantum predictions \\
1982         & Aspect, Grangier, Roger (experiment)
             & Landmark test: Bell inequalities violated, as quantum mechanics predicts \\
1996         & Rovelli (theoretical physics)
             & Relational QM: quantum states are always relative, never absolute \\
2022         & Zeilinger, Aspect, Clauser (Nobel Prize)
             & Nobel Prize for experiments establishing the violation of Bell inequalities \\
\bottomrule
\end{tabularx}
\end{table}

The \bahai writings were not responding to quantum mechanics --- they
preceded it by a century.
They were not responding to Whitehead --- they preceded him by decades.
They arrived where they arrived by their own means: spiritual discernment,
articulated in the writings of \Bahaullah and \AbdulBaha.

This is not presented as proof of supernatural authority.
What it establishes, on my reading, is independence: three traditions that did not communicate
with each other, using three different methods, arrived at what I read as the same structural
description of reality.
That convergence requires explanation.

A reader may reply that structural coincidences across traditions are not rare
in the history of ideas, and that no particular convergence demands a
particular metaphysics.
The reply is fair as a general caution.
What it does not address is the specificity of the present convergence.
Six independent structural points, drawn from texts that did not communicate
with each other, agreeing not in mood or imagery but in logical architecture:
this is a degree of correspondence for which the appeal to coincidence
becomes itself a load-bearing claim, and one that requires its own argument.
I do not insist on a metaphysical conclusion here.
I insist only that the convergence is real, that it is specific, and that
explanations of it must take its specificity into account.

% ── SECTION 5: QUANTUM PHYSICS AS CONVERGENCE, NOT APPROPRIATION (~500 words)

\section{Quantum Physics as Convergence, Not Appropriation}
\label{sec:confirmation}

% This project does not claim \Bahaullah was writing physics.
% Claim: the structural description in the writings CORRESPONDS to what
% physics discovered independently.
% Relationship: confirmation, not borrowing.
% Each tradition keeps its own vocabulary and methodology.

A clarification is essential.

I do not claim that \Bahaullah was writing physics, or that the
\bahai writings anticipate the technical formalism of quantum mechanics.
The writings are not a physics textbook.
They are spiritual revelation, addressed to questions of human meaning,
moral development, and the nature of God.

What I claim is different: that the \emph{structural description}
of reality in the \bahai writings --- existence as relational affinity, creation
as relational event --- corresponds to the relational structure that quantum
mechanics, on the relational reading, has independently arrived at through a
century of mathematics and experiment.
The third element of that description, love as the creative foundation at every
degree, has no counterpart in physics; it is where the writings, and Whitehead
in his own vocabulary, go further.

The relationship is convergence, not appropriation.
It is not confirmation either: experiments test physical theories, not the
writings, and the agreement lies in the shape of the two descriptions.
The \bahai writings do not need quantum mechanics to be true.
Quantum mechanics does not need the \bahai writings to be true.
What I document is that, independently of each other, both arrived at
the same relational structure --- which is what the two-wings principle predicts they should do.

Each tradition keeps its own vocabulary and methodology.
A \bahai scholar does not need to accept the formalism of the Relational Coherence
Framework to validate my argument.
A physicist does not need to accept the authority of the \bahai writings.
Each can assess the convergence on its own grounds.

% ── SECTION 6: THE NON-COMPOSITE SOUL AND DECOHERENCE (~500 words) ──────────

\section{The Non-Composite Soul: Implications}
\label{sec:soul}

% \bahai: soul is not a combination of elements (\AbdulBaha, SAQ).
% What this implies within a relational framework: no internal structure
% to lose coherence between.
% NOT a claim that the soul is a quantum system.
% A structural observation about what non-compositeness entails.

\AbdulBaha teaches that the soul ``is not a combination of elements,
it is not composed of many atoms''~\citep{SAQ}.
The soul is non-composite: it has no constituent parts, no internal structure
from which it could be assembled or into which it could be dissolved.

Within the relational ontological framework I have established in this paper,
this claim has a precise implication.
Quantum decoherence is a process by which a system's internal degrees of freedom
become entangled with environmental degrees of freedom, causing the system to
lose quantum coherence.
But decoherence requires \emph{internal structure}: there must be parts
between which coherence can be lost.

A genuinely non-composite entity has no internal degrees of freedom.
It therefore cannot undergo quantum decoherence in the technical sense.

This is not a claim that the soul is a quantum system, or that the soul exists
in superposition, or that it can be described by a wave function.
It is a structural observation: if the soul is what the \bahai writings say it is
--- non-composite, not assembled from elements --- then, on my reading, the standard decoherence
argument does not apply to it.

Two implications follow, and both are best stated as questions I open rather
than as positions I defend.

The first concerns the long-standing decoherence objection in the quantum
mind debate.
Critics of quantum-mechanical accounts of consciousness have argued that the
warm and noisy environment of the brain destroys quantum coherence too rapidly
for any cognitive process to depend on it~\citep{Zurek2003}.
The objection is decisive against any account that locates the relevant
coherence in macromolecular structures whose decoherence times are short.
It does not, however, apply with the same force to an entity that has no
internal structure at all.
Whether the \bahai conception of the soul can do work in this debate is a
question for further inquiry; what \S\ref{sec:soul} establishes is only that
the standard decoherence objection cannot be redirected against the
non-composite soul without modification, since the objection presupposes
exactly the internal compositeness the \bahai conception denies.

The second concerns the role of the observer.
Relational quantum mechanics makes the relational interaction between system
and observer constitutive of what becomes definite~\citep{Rovelli1996}.
If a non-composite, relationally embedded entity is what \bahai scholarship
has in mind by the soul, then, on my reading, such an entity is structurally available as an
observer in the precise relational sense the physics requires.
My argument does not establish that the soul plays this role.
It establishes only that the structural conditions for it to do so are not
incoherent.
A research programme oriented to this question would draw on \bahai
philosophical anthropology, on the relational reading of quantum mechanics,
and on the conservation arguments developed in \S\ref{sec:convergences}.
I mark the programme as open and leave its substantive investigation to
subsequent work.

% ── SECTION 7: CONCLUSION (~600 words) ───────────────────────────────────────

\section{Conclusion: The \bahai Contribution to Science-Religion Scholarship}
\label{sec:conclusion}

% Key points:
%   - The \bahai writings contribute something unique: a spiritually revealed
%     account that arrived, by a different method, at the same structural
%     conclusions as modern physics and process philosophy.
%   - This is the two-wings principle — not as hope or doctrine, but as fact.
%   - Opens new channels for academic \bahai-science dialogue.

The \bahai writings contribute something unique to the science-religion
conversation.
They constitute a spiritually revealed account of reality that arrived,
by a different method, at the same structural conclusions as modern physics
and process philosophy --- forty years before Whitehead named the framework
and a century before physics confirmed it.

This is the two-wings principle --- not as a hope, not as a doctrine, but as
a demonstrated finding.
On my reading, science and religion have, independently of each other, arrived at the same
description of what reality is made of: relationship, constituted by relational
affinity, sustained by what each tradition calls love.

The implications for \bahai scholarship are significant.
The principle of the harmony of science and religion can now be defended not
merely by assertion but by argument --- by showing, in the language of modern
physics and philosophy, that the structural architecture of the \bahai
metaphysics of \mahabbat, on my reading, matches the structural architecture that
physics has empirically confirmed at the most fundamental level of physical
description we possess.

That is not a small finding.
It is the two wings, confirmed.
\newpage

\section*{Declaration of AI-Assisted Technologies}
\addcontentsline{toc}{section}{Declaration of AI-Assisted Technologies}

This manuscript was prepared with the assistance of Claude (Anthropic), a large language
model, used in the capacities detailed here for transparency, including a set of late
revisions made in September 2026.

\textbf{Argument and analysis.} The central textual and historical argument of this
paper --- the chronological case for the priority of the Bahá'í writings' relational
metaphysics of \emph{mahabbat} (Section~2), the six structural convergences with quantum
mechanics and process philosophy (Section~3), the chronological framing (Section~4), and
the treatment of the non-composite soul (Section~6, except the revised decoherence
argument in Section~6.1 described below) --- is my own analysis. Claude was used in an editorial and research-assistant
capacity for these sections: expanding and tightening prose I had outlined, and verifying
citations against primary sources.

\textbf{AI-drafted framing sections.} The introduction (Section~1), the discussion
distinguishing convergence from appropriation (Section~5), and the conclusion
(Section~7) were drafted by Claude from my direction and outline, then reviewed and
revised by me across multiple passes before inclusion.

\textbf{Citation and reference verification.} Claude fetched primary-source texts
directly (from \url{www.bahai.org} and \url{reference.bahai.org}) and checked the
manuscript's verbatim quotations and load-bearing citations against them; I reviewed and
approved each correction that resulted. While not framed as an exhaustive line-by-line
audit of every paraphrase citation, this check covered every direct quotation and every
citation whose accuracy the argument depends on --- including a final pass over the two
paraphrase citations previously left open, both now resolved (one verified correct
against \emph{Some Answered Questions} Part~4, one corrected in \emph{Selections from the
Writings of `Abdu'l-Bah\'a}) --- and it identified and corrected several
errors introduced during drafting: a citation to \emph{Some Answered Questions} for a
quotation that in fact originates in \emph{Paris Talks}; two passages presented as
verbatim quotations from \emph{The Promulgation of Universal Peace} whose closing
clauses did not match the source text and have been corrected to the verified wording;
and a placeholder bibliography entry (an unverifiable single-volume anthology) that has
been replaced with the two specific, citable primary sources the quotations actually
come from.

\textbf{Late revisions (September 2026).} After the manuscript was first prepared for
submission, I had Claude review it against my companion book manuscript and against the
physics; the resulting revisions were made at my direction on 29 September 2026. Claude
identified that the original argument of Section~6.1 --- that decoherence requires
internal parts, so an entity without parts cannot decohere --- is refuted by a single
electron, which has no known parts yet decoheres readily, and drafted the revised
argument, which rests the soul's immunity to decoherence on its not being a physical
system and states the interaction problem as open; Section~6.1 as it now stands is
therefore substantially AI-drafted, from my direction. Claude also revised the abstract,
introduction, Sections~2--5 and the conclusion so that the paper claims structural
convergence rather than confirmation and attributes the claim that love sustains
relationship to the Bahá'í writings (and to Whitehead) rather than to physics; corrected
overstatements about the holographic principle (Section~3.5) and about zero-point energy
and the quantum vacuum (Sections~2 and~3.6); and re-checked each direct quotation against
the specific work, and for \emph{The Promulgation of Universal Peace} the specific talk,
that it cites, restoring one missing ellipsis.

% TODO(author): before submitting, review the September 2026 revisions (especially
% Section 6.1) so that the responsibility statement below is accurate for them too.
I have reviewed all AI-assisted and AI-drafted content, take full responsibility for the
accuracy and integrity of the manuscript's claims and citations, and affirm that the
argument, its interpretation, and its conclusions are my own.

\section*{Works Cited}
\bibliography{references}

\end{document}
